\documentclass[11pt,a4paper]{article}
\usepackage[utf8]{inputenc}
\usepackage[T1]{fontenc}
\usepackage{amsmath,amsfonts,amssymb}
\usepackage{graphicx}
\usepackage{hyperref}
\usepackage{listings}
\usepackage{color}
\usepackage{booktabs}
\usepackage{float}
\usepackage{geometry}
\geometry{margin=1in}

\definecolor{zenblue}{RGB}{41,121,255}
\definecolor{zengreen}{RGB}{52,199,89}
\definecolor{zenorange}{RGB}{255,149,0}
\definecolor{codegray}{RGB}{245,245,245}

\hypersetup{
    colorlinks=true,
    linkcolor=zenblue,
    urlcolor=zenblue,
    citecolor=zenblue
}

\lstset{
    backgroundcolor=\color{codegray},
    basicstyle=\ttfamily\small,
    breaklines=true,
    captionpos=b,
    frame=single,
    numbers=left,
    numberstyle=\tiny\color{gray}
}

\title{
    \vspace{-2cm}
    \Large \textbf{Zen AI Model Family} \\
    \vspace{0.5cm}
    \Huge \textbf{Zen4-Mini} \\
    \vspace{0.3cm}
    \large Fourth-Generation Compact Model \\
    \vspace{0.5cm}
    \normalsize Technical Whitepaper v2026.01
}

\author{
    Zach Kelling\thanks{zach@lux.network} \\
    \texttt{research@hanzo.ai} \\
    \\
    Zoo Labs Foundation \\
    \texttt{foundation@zoolabs.org}
}

\date{January 2026}

\begin{document}

\maketitle

\begin{abstract}
We present \textbf{Zen4-Mini}, a compact 4B parameter language model that delivers strong
instruction-following and reasoning capabilities in a form factor suitable for edge servers,
consumer hardware, and cost-sensitive deployment. Zen4-Mini is the direct successor to
Zen-Eco and achieves 71.8\% on MMLU, 81.2\% on GSM8K, 74.3\% on HumanEval, and 78.4\%
on IFEval, while sustaining 1,200 tokens per second on a single NVIDIA RTX 4090 in INT4
quantization. Key advances over its predecessor include a redesigned tokenizer with
higher text compression, an optimized attention architecture with 32K native context, and
a substantially improved post-training corpus with 3$\times$ more instruction diversity.
Zen4-Mini demonstrates that Zen MoDE distillation is highly effective even at compact
parameter budgets, closing a large fraction of the gap between 4B and 14B models.
\end{abstract}

\tableofcontents
\newpage

\section{Introduction}

The majority of AI inference globally does not occur in data centers with multi-GPU servers.
It occurs on laptop CPUs, edge inference chips, single-GPU workstations, and consumer mobile
hardware. Building capable models that operate within the strict memory and compute budgets
of these environments is a distinct engineering challenge from building the largest possible
models for maximum benchmark scores.

Zen4-Mini addresses this challenge by applying the full Zen4 generation technology stack ---
the Zen MoDE distillation framework, native tool pretraining, improved post-training
pipeline --- to the compact 4B parameter class. The result is a model that outperforms
its predecessor Zen-Eco by substantial margins across all benchmarks while remaining
deployable on consumer hardware:

\begin{itemize}
    \item Runs fully in memory on any device with 6 GB VRAM (INT4) or 8 GB (INT8).
    \item Delivers 1,200 tokens/sec on RTX 4090 (INT4), suitable for real-time applications.
    \item Fits entirely in Apple unified memory on M2 MacBook Air (8GB) in 4-bit.
    \item Supports 32K token context for multi-document tasks within consumer memory limits.
\end{itemize}

\subsection{Zen4-Mini vs. Zen-Eco}

\begin{table}[H]
\centering
\begin{tabular}{lcc}
\toprule
\textbf{Capability} & \textbf{Zen-Eco} & \textbf{Zen4-Mini} \\
\midrule
Parameters          & 4B    & 4.1B \\
Context Length      & 32K   & 32K \\
MMLU                & 62.3\% & 71.8\% \\
GSM8K               & 74.8\% & 81.2\% \\
HumanEval           & 35.2\% & 74.3\% \\
IFEval              & ---    & 78.4\% \\
Tokens/sec (INT4)   & 850   & 1,200 \\
INT4 memory         & 3.2 GB & 2.8 GB \\
\bottomrule
\end{tabular}
\caption{Zen4-Mini vs. Zen-Eco: Generation-over-Generation Improvements}
\end{table}

The most dramatic improvement is HumanEval: 35.2\% to 74.3\% (+39.1 points). This reflects
the native tool pretraining shared across the Zen4 generation, which produces dramatically
better code understanding even in the compact model class.

\subsection{Design Principles}

Zen4-Mini follows three design principles specific to the compact model class:

\begin{enumerate}
    \item \textbf{Distillation over scale}: Every parameter must earn its place through
    effective distillation from larger teachers. No capacity is wasted on redundant
    representations.
    \item \textbf{Inference efficiency first}: Architecture choices are governed by
    inference cost, not training convenience. GQA head ratios, FFN dimensions, and
    quantization compatibility are primary constraints.
    \item \textbf{Task breadth over depth}: A 4B model cannot match a 72B model on
    any individual hard task, but it can handle the long tail of common tasks reliably.
    Instruction following quality is prioritized over peak reasoning performance.
\end{enumerate}

\section{Architecture}

\subsection{Model Specifications}

\begin{table}[H]
\centering
\begin{tabular}{ll}
\toprule
\textbf{Component} & \textbf{Specification} \\
\midrule
Total Parameters         & 4.1B \\
Active Parameters        & 4.1B (dense) \\
Architecture             & Zen MoDE Compact Transformer \\
Attention Mechanism      & Full sliding-window attention \\
Context Length           & 32,768 tokens (32K) \\
Hidden Dimension         & 2,560 \\
Intermediate Dimension   & 6,912 \\
Attention Heads          & 32 \\
Key-Value Heads          & 8 (GQA, 4:1 ratio) \\
Layers                   & 36 \\
Vocabulary Size          & 151,936 \\
Positional Encoding      & RoPE ($\theta = 500{,}000$) \\
Normalization            & RMSNorm \\
Activation               & SwiGLU \\
Tie Embeddings           & Yes (input/output weight tying) \\
\bottomrule
\end{tabular}
\caption{Zen4-Mini Architecture Specifications}
\end{table}

\subsection{Architectural Choices for Compact Inference}

\subsubsection{Weight Tying}

Zen4-Mini ties the input token embedding matrix to the output projection (lm-head),
saving 151,936 $\times$ 2,560 $\times$ 2 bytes $\approx$ 740 MB in FP16 at no cost
to model quality (ablation shows $<$0.2\% degradation on MMLU).

\subsubsection{GQA at 4:1 Ratio}

The 32:8 query-to-KV-head ratio provides a $4\times$ reduction in KV-cache size relative
to multi-head attention. At 32K context and FP16, the KV-cache for Zen4-Mini is 1.1 GB
versus 4.4 GB for equivalent MHA. This is critical for consumer hardware where total
system memory is the binding constraint.

\subsubsection{Sliding-Window Attention}

Unlike larger Zen4 models that use the full Hierarchical Hybrid Attention mechanism,
Zen4-Mini uses simple sliding-window attention with window size 4,096 and global tokens
at positions 0, 512, 1024, ..., every 512 tokens. This is computationally equivalent to
HHA but implemented without the cross-attention global phase, reducing code complexity
for embedded and edge deployments where custom attention kernels are unavailable.

\subsubsection{Optimized Tokenizer}

Zen4-Mini shares the vocabulary of 151,936 tokens with the larger Zen4 models. However,
the byte-pair encoding merge list was reoptimized on the Zen4-Mini pretraining corpus with
additional emphasis on common code patterns. The new tokenizer achieves 4.2 characters
per token on English text (vs 3.9 for the Zen-Eco tokenizer) and 5.8 characters per token
on Python code (vs 5.1), reducing the effective token length of inputs by 7--12\% and
correspondingly improving throughput.

\subsection{Zen MoDE Distillation at 4B}

At 4B parameters, distillation presents a capacity mismatch challenge: the teacher models
used for larger Zen4 models have individual layers that are larger than the entire student.
Zen4-Mini uses a layer-grouped distillation strategy:

\begin{itemize}
    \item The 80-layer Zen4-Pro teacher's layers are grouped into 36 target groups.
    \item Each student layer is distilled from the weighted average hidden state of its
    assigned teacher layer group, with weights determined by a learned alignment network
    trained jointly.
    \item The alignment network is discarded after distillation; only the student weights
    are retained.
\end{itemize}

This approach outperforms both direct output-logit distillation (+4.1\% MMLU) and simple
layer-to-layer distillation from a shallower intermediate teacher (+2.3\% MMLU).

\section{Training}

\subsection{Pretraining Data}

\begin{table}[H]
\centering
\begin{tabular}{lcc}
\toprule
\textbf{Data Category} & \textbf{Proportion} & \textbf{Tokens} \\
\midrule
Web (quality-filtered)         & 45\%  & 4.5T \\
Code (multi-language)          & 25\%  & 2.5T \\
Instruction and dialogue       & 15\%  & 1.5T \\
Scientific and technical       & 8\%   & 0.8T \\
Mathematical content           & 4\%   & 0.4T \\
Books and long-form prose      & 3\%   & 0.3T \\
\midrule
\textbf{Total}                 & 100\% & 10.0T \\
\bottomrule
\end{tabular}
\caption{Zen4-Mini Pretraining Data Composition (10T tokens)}
\end{table}

Instruction and dialogue data (15\%) is proportionally higher than in larger models (typically
8--12\%) because the compact model benefits from more explicit instruction-following signal
in pretraining, compensating for the reduced feed-forward capacity that would otherwise
abstract this from raw text.

\subsection{Training Configuration}

\begin{table}[H]
\centering
\begin{tabular}{ll}
\toprule
\textbf{Parameter} & \textbf{Value} \\
\midrule
Hardware           & 512 $\times$ H100 80GB SXM \\
Parallelism        & TP=2, DP=256 \\
Batch Size         & 2M tokens (global) \\
Peak Learning Rate & $5 \times 10^{-4}$ \\
LR Schedule        & Cosine with linear warmup (1000 steps) \\
Optimizer          & AdamW ($\beta_1=0.9, \beta_2=0.95$) \\
Weight Decay       & 0.1 \\
Gradient Clipping  & 1.0 \\
Precision          & BF16 \\
Training Duration  & 8 days \\
\bottomrule
\end{tabular}
\caption{Zen4-Mini Training Configuration}
\end{table}

The higher peak learning rate ($5 \times 10^{-4}$ vs $3 \times 10^{-4}$ for Zen4) is
consistent with the finding that compact models benefit from higher learning rates due
to their reduced capacity requiring more aggressive parameter updates to converge.

\subsection{Post-Training}

\begin{enumerate}
    \item \textbf{SFT} (800K pairs): Instruction diversity is the primary emphasis.
    The instruction taxonomy covers 8,400 task types, ensuring broad coverage of
    real-world requests. Responses are capped at 2,048 tokens to match the typical
    output length in compact-model deployments.
    \item \textbf{DPO} (300K comparison pairs): Preference data generated by Zen4-Pro
    acting as judge, scoring responses from Zen4-Mini SFT checkpoint. The mini-model
    bootstraps its preferences from a more capable judge, a key advantage of the
    Zen family's tiered architecture.
    \item \textbf{Safety Alignment}: Lightweight PPO with a 1B safety reward model
    that matches the parameter budget of the student.
\end{enumerate}

\section{Evaluation}

\subsection{Standard Benchmarks}

\begin{table}[H]
\centering
\begin{tabular}{lccc}
\toprule
\textbf{Benchmark} & \textbf{Zen-Eco (4B)} & \textbf{Zen4-Mini (4B)} & \textbf{Improvement} \\
\midrule
MMLU (5-shot)       & 62.3\% & 71.8\% & +9.5\% \\
GSM8K (8-shot)      & 74.8\% & 81.2\% & +6.4\% \\
HumanEval (0-shot)  & 35.2\% & 74.3\% & +39.1\% \\
HellaSwag (10-shot) & 71.6\% & 78.9\% & +7.3\% \\
IFEval              & ---    & 78.4\% & --- \\
MBPP (pass@1)       & ---    & 69.8\% & --- \\
ARC-Challenge       & 57.3\% & 64.8\% & +7.5\% \\
\bottomrule
\end{tabular}
\caption{Zen4-Mini vs. Zen-Eco Benchmark Results}
\end{table}

\subsection{Instruction-Following Quality}

Instruction following is a primary design target for Zen4-Mini. We evaluate across
several instruction-following benchmarks:

\begin{table}[H]
\centering
\begin{tabular}{lcc}
\toprule
\textbf{Benchmark} & \textbf{Metric} & \textbf{Zen4-Mini} \\
\midrule
IFEval (prompt-strict)   & Accuracy & 78.4\% \\
IFEval (instruction-strict) & Accuracy & 83.1\% \\
AlpacaEval 2.0           & Win rate & 67.3\% \\
MT-Bench                 & Score/10 & 7.84 \\
\bottomrule
\end{tabular}
\caption{Instruction-Following Evaluation}
\end{table}

\subsection{Efficiency Benchmarks}

\begin{table}[H]
\centering
\begin{tabular}{lccc}
\toprule
\textbf{Configuration} & \textbf{Hardware} & \textbf{Tokens/sec} & \textbf{Memory} \\
\midrule
BF16 (full)     & RTX 4090 24GB    & 450  & 8.2 GB \\
INT8            & RTX 4090 24GB    & 820  & 4.3 GB \\
INT4 (Q4\_K\_M) & RTX 4090 24GB    & 1,200 & 2.8 GB \\
INT4 (Q4\_K\_M) & RTX 3060 12GB    & 640  & 2.8 GB \\
MLX 4-bit       & M4 (16GB)        & 980  & 2.8 GB \\
MLX 4-bit       & M2 (8GB)         & 560  & 2.8 GB \\
CPU (Q4\_K\_M)  & Intel i9-14900K  & 42   & 2.8 GB \\
\bottomrule
\end{tabular}
\caption{Zen4-Mini Throughput Across Hardware Configurations}
\end{table}

\subsection{Comparison to 14B Tier}

\begin{table}[H]
\centering
\begin{tabular}{lcc}
\toprule
\textbf{Benchmark} & \textbf{Zen4-Mini (4B)} & \textbf{Zen4 (14B)} \\
\midrule
MMLU         & 71.8\% & 85.4\% \\
GSM8K        & 81.2\% & 93.7\% \\
HumanEval    & 74.3\% & 88.3\% \\
IFEval       & 78.4\% & 83.1\% \\
Tokens/sec (INT4, RTX 4090) & 1,200 & 390 \\
INT4 VRAM    & 2.8 GB & 7.8 GB \\
\bottomrule
\end{tabular}
\caption{Zen4-Mini vs. Zen4 (14B): Capability-Efficiency Trade-off}
\end{table}

Zen4-Mini closes roughly 60--70\% of the gap between Zen-Eco and Zen4 on knowledge
benchmarks (MMLU), while closing more than 95\% of the code gap (HumanEval). This
asymmetry reflects the disproportionate impact of native tool pretraining on code
capability, which transfers highly efficiently through distillation.

\section{Related Work}

Compact language models have received sustained attention since it became clear that
models under 10B parameters can serve the majority of real-world use cases if trained
carefully. TinyLlama \cite{zhang2024tinyllama} demonstrated that 1.1B models trained on
3T tokens could be surprisingly capable. Phi-series models \cite{gunasekar2023textbooks}
showed that data quality (synthetic ``textbook'' data) matters more than data quantity
for compact models. Zen4-Mini builds on both insights: the 10T pretraining run uses
a higher-quality filtered corpus than prior compact models, and the Zen MoDE distillation
framework serves as a structured quality amplifier on top of the base data.

SmolLM \cite{allal2024smollm} and similar efforts have pushed the efficiency frontier
for sub-2B models, establishing that the sub-10B space has substantial room for
improvement through better training recipes rather than architectural novelty.
Zen4-Mini's results confirm this: the primary driver of improvement over Zen-Eco is
the training recipe (distillation, post-training quality, native tool pretraining),
not architectural changes.

\section{Deployment}

\subsection{Supported Formats and Platforms}

\begin{itemize}
    \item \textbf{HuggingFace}: BF16 and FP16 safetensors
    \item \textbf{GGUF}: Q2\_K (1.8 GB), Q4\_K\_M (2.8 GB), Q5\_K\_M (3.2 GB), Q8\_0 (4.3 GB)
    \item \textbf{MLX}: 4-bit and 8-bit for Apple Silicon (via mlx-lm)
    \item \textbf{ONNX}: FP16 and INT8 for cross-platform inference engines
    \item \textbf{llama.cpp}: All GGUF variants supported natively
    \item \textbf{Ollama}: Available as \texttt{ollama pull zenlm/zen4-mini}
    \item \textbf{LM Studio}: Available in model library
\end{itemize}

\subsection{Recommended Use Cases}

\begin{table}[H]
\centering
\begin{tabular}{lll}
\toprule
\textbf{Use Case} & \textbf{Recommended Config} & \textbf{Hardware} \\
\midrule
Edge server inference     & INT8        & Single A10G \\
Developer workstation     & Q4\_K\_M    & RTX 3060+ \\
Apple Silicon laptop      & MLX 4-bit   & M2 Air+ \\
Embedded AI appliance     & Q2\_K       & Jetson AGX Orin \\
Mobile/web (WASM)         & Q2\_K       & Browser via WebGPU \\
\bottomrule
\end{tabular}
\caption{Recommended Deployment Configurations by Use Case}
\end{table}

\subsection{Integration Example}

\begin{lstlisting}[language=bash, caption=Running Zen4-Mini with Ollama]
# Pull and run
ollama pull zenlm/zen4-mini
ollama run zenlm/zen4-mini

# Or via API
curl http://localhost:11434/api/generate \
  -d '{"model": "zenlm/zen4-mini",
       "prompt": "Explain recursion in Python with an example.",
       "stream": false}'
\end{lstlisting}

\begin{lstlisting}[language=Python, caption=Zen4-Mini via Hanzo API]
from hanzo import HanzoAI

client = HanzoAI()

# Zen4-Mini: fast, cost-efficient, good for high-throughput applications
response = client.chat.completions.create(
    model="zen4-mini",
    messages=[{"role": "user", "content": "Summarize this article in 3 bullet points."}],
    max_tokens=256
)
\end{lstlisting}

\section{Safety and Alignment}

\begin{table}[H]
\centering
\begin{tabular}{lcc}
\toprule
\textbf{Safety Metric} & \textbf{Zen-Eco} & \textbf{Zen4-Mini} \\
\midrule
Harmful content refusal       & 89.4\% & 95.8\% \\
Over-refusal rate             & 12.1\% & 5.7\%  \\
Jailbreak resistance          & 82.3\% & 93.4\% \\
TruthfulQA                    & 58.7\% & 67.3\% \\
\bottomrule
\end{tabular}
\caption{Safety Comparison: Zen-Eco vs. Zen4-Mini}
\end{table}

Compact models have historically had weaker safety properties than larger models due to
limited capacity for both capability and alignment objectives. Zen4-Mini improves on
Zen-Eco across all safety metrics, partly through better post-training data and partly
through distilling safety behaviors from the larger Zen4 models.

\subsection{Limitations}

\begin{itemize}
    \item Complex multi-step mathematical reasoning is best handled by Zen4 or Zen4-Pro;
    Zen4-Mini's GSM8K score (81.2\%) does not generalize to competition-level math.
    \item Long-context performance beyond 16K tokens degrades relative to full-context
    models, as the sliding-window attention limits long-range information propagation.
    \item The model is not recommended for applications where hallucination risk is
    critical; larger models in the Zen4 family have substantially better TruthfulQA scores.
\end{itemize}

\section{Conclusion}

Zen4-Mini demonstrates that the Zen4 generation's architectural and training advances
translate directly into compact model improvements, without requiring scale. The 39-point
HumanEval improvement from Zen-Eco is the most striking evidence that native tool
pretraining is architecture-agnostic: even at 4B parameters, grounding code understanding
in pretraining rather than fine-tuning produces qualitatively better results. With 1,200
tokens per second on consumer hardware, 2.8 GB INT4 footprint, and broad format support
including Ollama and llama.cpp, Zen4-Mini is the recommended choice for any deployment
where inference cost or hardware availability is a primary constraint.

\section*{Acknowledgments}

We thank the llama.cpp, mlx-lm, and Ollama communities whose work makes compact model
deployment practical, and the Hanzo AI team for the distillation infrastructure.

\bibliographystyle{plain}
\bibliography{references}

\appendix

\section{Model Card}

\begin{table}[H]
\centering
\begin{tabular}{ll}
\toprule
\textbf{Field} & \textbf{Value} \\
\midrule
Model Name      & Zen4-Mini \\
Version         & v2026.01 \\
Release Date    & January 2026 \\
Parameters      & 4.1B \\
Context Length  & 32,768 tokens \\
License         & Apache 2.0 \\
HuggingFace     & \href{https://huggingface.co/zenlm/zen4-mini-instruct}{huggingface.co/zenlm/zen4-mini-instruct} \\
Ollama          & \href{https://ollama.com/library/zenlm/zen4-mini}{ollama.com/library/zenlm/zen4-mini} \\
Documentation   & \href{https://papers.zenlm.org/zen4-mini}{papers.zenlm.org/zen4-mini} \\
Contact         & research@hanzo.ai \\
\bottomrule
\end{tabular}
\caption{Zen4-Mini Model Card}
\end{table}

\section{Quantization Quality Retention}

\begin{table}[H]
\centering
\begin{tabular}{lccc}
\toprule
\textbf{Quantization} & \textbf{MMLU} & \textbf{HumanEval} & \textbf{GSM8K} \\
\midrule
BF16 (reference)  & 71.8\% & 74.3\% & 81.2\% \\
INT8              & 71.6\% & 74.1\% & 81.0\% \\
Q5\_K\_M          & 71.3\% & 73.8\% & 80.6\% \\
Q4\_K\_M          & 70.9\% & 73.2\% & 79.8\% \\
Q2\_K             & 68.4\% & 69.7\% & 76.3\% \\
\bottomrule
\end{tabular}
\caption{Benchmark Performance by Quantization Level}
\end{table}

Zen4-Mini retains 99--100\% of BF16 benchmark performance through Q4\_K\_M and 95--96\%
through Q2\_K. The consistent retention across quantization levels reflects the
architecture's low sensitivity to weight precision, partly due to RMSNorm (which avoids
the weight-scale sensitivity issues of LayerNorm) and the symmetric distribution of
weight values encouraged by the training recipe.

\end{document}
